\section{Entire actual diagonal preimages}
\begin{lemma}[The same mixed map has the actual original diagonal preimages]
\label{lem:central-embedding-same-mixed-map-entire-diagonal-preimages}
\uses{lem:central-embedding-same-mixed-projective-map-entire-chart,def:central-embedding-actual-mixed-padded-pair-coordinate}
\lean{EnvelopeProperty.ProjectiveCentralFiberEmbeddingNeighborhoodClosedMixedCoordinatePreimages.mixedProjectiveMap_preimage_diagonal}
\leanok
The WHOLE inverse image of the actual $(i,i)$ target coordinate open under
the SAME native mixed generated map is exactly the original $i$-th coordinate
open. No preimage, family, map or chart matching certificate is supplied.
The domain assumption is only on the ambient coefficient ring; it adds no
integral or reduced hypothesis to a later actual central fibre.
\end{lemma}
\begin{proof}
Use the actual original source coordinate cover, so every original point lies
in an actual polynomial source chart $k$. The proved entire chart formula
identifies the SAME generated morphism there with the actual native mixed
tuple map. The native unit-ideal tuple preimage theorem identifies the entire
inverse target coordinate open with the original spectrum basic open of that
actual tuple entry. The native finite equivalence and original zero-padding
convention identify the $(i,i)$ entry with
$(X_i/X_k)^d(X_i/X_k)=(X_i/X_k)^{d+1}$. Its actual spectrum basic open is the
basic open of $X_i/X_k$, since $d+1$ is positive. The native original coordinate
chart preimage theorem identifies this with the entire pullback of the
original $i$-th source coordinate open. This proves the equality on the actual
cover and hence on the whole original scheme, including nonclosed points.
\end{proof}
The diagonal target opens are not asserted to cover the entire target. They
cover the actual image after this whole-preimage identity. Actual local chart
closed immersions and properness must still give global closedness; neither
is an input of this theorem.
